\begin{tikzpicture}[
  font=\scriptsize,
  panah/.style={-{Stealth[length=2mm]}, thick},
  persp/.style={draw, rounded corners, fill=blue!8, minimum width=2cm,
    minimum height=0.6cm, align=center},
  ranah/.style={draw, rounded corners, fill=orange!12, minimum width=2.2cm,
    minimum height=0.8cm, align=center},
  fase/.style={draw, rounded corners, fill=blue!8, minimum width=2.4cm,
    minimum height=0.9cm, align=center}
]
  % Enam perspektif, dikelompokkan menjadi kapabilitas bisnis dan teknis.
  \node[persp] (bus) at (0.9,2.5) {Business};
  \node[persp] (peo) at (0.9,1.8) {People};
  \node[persp] (gov) at (0.9,1.1) {Governance};
  \node[persp] (pla) at (0.9,-0.4) {Platform};
  \node[persp] (sec) at (0.9,-1.1) {Security};
  \node[persp] (ops) at (0.9,-1.8) {Operations};
  \node[draw, dashed, gray, rounded corners, inner sep=4pt,
    fit=(bus)(gov)] (bisnis) {};
  \node[draw, dashed, gray, rounded corners, inner sep=4pt,
    fit=(pla)(ops)] (teknis) {};
  \node[above=0pt of bisnis, gray!60!black] {{\tiny\mdseries kapabilitas bisnis}};
  \node[above=0pt of teknis, gray!60!black] {{\tiny\mdseries kapabilitas teknis}};

  % Capability dasar menjadi penghubung perspektif dan domain transformasi.
  \node[draw, rounded corners, fill=praditagreen!15, minimum width=2.2cm,
    minimum height=5.1cm, align=center, text width=1.9cm,
    font=\scriptsize\bfseries] (cap) at (4.2,0.35)
    {\textit{Capability}\\dasar\\[3pt]{\tiny\mdseries dimiliki atau dikelola
     pemangku kepentingan tiap perspektif\par}};

  % Empat domain transformasi.
  \node[ranah] (tek) at (7.4,2.15) {Technology};
  \node[ranah] (pro) at (7.4,0.85) {Process};
  \node[ranah] (org) at (7.4,-0.45) {Organization};
  \node[ranah] (prd) at (7.4,-1.75) {Product};
  \node[above=2pt of tek, gray!60!black] {{\tiny\mdseries domain transformasi}};

  % Hasil bisnis sebagai tujuan akhir.
  \node[draw, rounded corners, fill=green!10, minimum width=2.2cm,
    minimum height=2.2cm, align=center, text width=1.9cm,
    font=\scriptsize\bfseries] (hasil) at (10.6,0.2)
    {Hasil bisnis\\[3pt]{\tiny\mdseries risiko turun, pendapatan tumbuh, operasi lebih
     efisien\par}};

  \draw[panah] (bisnis.east) -- (bisnis.east -| cap.west);
  \draw[panah] (teknis.east) -- (teknis.east -| cap.west);
  \foreach \ranah in {tek, pro, org, prd} {
    \draw[panah] (cap.east |- \ranah) -- (\ranah.west);
    \draw[thick] (\ranah.east) -- ++(0.4,0);
  }
  \draw[thick] (8.9,2.15) -- (8.9,-1.75);
  \draw[panah] (8.9,0.2) -- (hasil.west);

  % Empat fase perjalanan transformasi, diulang per inisiatif.
  \node[gray!60!black] at (5.95,-2.65)
    {{\tiny\mdseries empat fase perjalanan transformasi}};
  \node[fase] (env) at (1.3,-3.3) {Envision\\{\tiny\mdseries peluang dan hasil}};
  \node[fase] (ali) at (4.4,-3.3) {Align\\{\tiny\mdseries celah kesiapan}};
  \node[fase] (lau) at (7.5,-3.3) {Launch\\{\tiny\mdseries pilot bernilai}};
  \node[fase, fill=green!10] (sca) at (10.6,-3.3)
    {Scale\\{\tiny\mdseries perluasan bertahap}};
  \draw[panah] (env) -- (ali);
  \draw[panah] (ali) -- (lau);
  \draw[panah] (lau) -- (sca);
  \draw[panah, gray] (sca.south) -- ++(0,-0.35) -| (env.south);
  \node[gray!60!black] at (5.95,-4.35)
    {{\tiny\mdseries diulang untuk setiap inisiatif berikutnya}};
\end{tikzpicture}
